\documentclass[11pt,a4paper]{article}

\usepackage[utf8]{inputenc}
\usepackage[T1]{fontenc}
\usepackage{lmodern}
\usepackage[margin=1in]{geometry}
\usepackage{hyperref}
\usepackage{booktabs}
\usepackage{enumitem}
\usepackage{amsmath}
\usepackage{graphicx}
\usepackage{xcolor}

\hypersetup{colorlinks=true, linkcolor=blue!60!black, urlcolor=blue!60!black, citecolor=blue!60!black}
\setlength{\parskip}{0.5em}
\setlength{\parindent}{0pt}
\pagestyle{plain}

\title{\Large\bfseries Fate-Engineered Oxidative Hair Dye Chemistry:\\
Precursors and Chromophores Designed to Be Stable in Hair,\\
Cleavable on Skin, and Degradable in Wastewater}
\author{Research Proposal (no experimental results reported) \\ Domain: Cosmetic Chemistry, Green Chemistry, Fiber Science, Toxicology}
\date{\today}

\begin{document}
\maketitle

\begin{abstract}
Permanent oxidative hair dyes depend on small aromatic amines that penetrate the cortex and couple into wash-fast chromophores. Their hazards are skin sensitization and persistent effluent. Existing improvements address one problem at a time: ME-PPD lowers sensitization, amino-acid alkalizers replace ammonia, and melanin-mimetic coatings avoid aromatic amines but lack shade range. We propose fate engineering: designing the precursor and the coupled dye so each behaves differently in three environments. In hair, the molecule persists and gives 28-shampoo fastness. On skin, an esterase-labile group converts unreacted precursor into a polar, less haptenating species. In wastewater, a trigger link breaks the chromophore for biodegradation. We describe a screening funnel of 5--8 cleavable PPD analogues against ME-PPD, with in-vitro sensitization assays, coupling kinetics, tress performance and OECD 301B testing. Arginine-based alkali and ME-PPD are used as the benchmark, not as the contribution. This is a research proposal; no experimental results are reported.
\end{abstract}

\noindent\textbf{Keywords:} oxidative hair dye, PPD analogues, ME-PPD, fate engineering, esterase-labile, wash-fastness, OECD 301B, sensitization

\tableofcontents
\newpage

\section{Introduction}

\subsection{Why oxidative dyes are hard to replace}
Over 50\% of women and a growing share of men in Europe, North America, the Middle East, and Asia use oxidative hair color. The permanent segment dominates value because only it covers grey durably in a single step.

Its mechanism imposes three linked requirements:
\begin{enumerate}[leftmargin=1.5em]
\item \textbf{Swell and open:} At pH 9--11, cuticle scales lift and the cortex swells so monomers can enter.
\item \textbf{Penetrate small, trap large:} Monomeric precursors (MW $\sim$100--150 Da) diffuse into the cortex. Peroxide oxidizes them with couplers into dimers and polymers (MW $>$ 500 Da) too large to diffuse out. This size-exclusion trap is permanence.
\item \textbf{Couple to color:} Primary intermediate + coupler $\rightarrow$ indo dye chromophore. PPD/PTD-based couplings remain unmatched for shade gamut and wash/light fastness.
\end{enumerate}

Green constraints pull against each step: lower pH reduces swelling and penetration; replacing small aromatic amines removes in-situ polymerization; weaker oxidation reduces coupling and lift; and a dye designed to biodegrade appears to contradict wash-fastness.

\subsection{Scope of this proposal}
This proposal does not claim a complete green permanent system, a new alkali, a new oxidant, or a full shade chart. Its scope is narrower: the design of the precursor and the coupled dye so that conditional stability resolves the fastness-versus-fate trade-off. Alkali and oxidant optimization, clinical testing, and salon trials are explicitly out of scope (see Section~\ref{sec:plan}).

\section{Prior art and gap}

\begin{table}[h]
\centering
\caption{Prior art and what it does not cover}
\label{tab:priorart}
\begin{tabular}{@{}lll@{}}
\toprule
Area & Prior art & Not covered \\
\midrule
Amino-acid alkali & Arginine/lysine alkalizer; & Permanent performance, \\
& arginine demi-permanents & effluent fate \\
Low-sensitizing precursor & ME-PPD; commercial use & Cross-elicitation; \\
& & no degradation design \\
Resorcinol-free & Resorcinol-free patents & Precursor persistence \\
Melanin-mimetic & DHI/DHICA; polydopamine & Narrow gamut; \\
& & surface coatings \\
Encapsulated oxidant & Encapsulated peroxide & Not tied to dye design \\
\bottomrule
\end{tabular}
\end{table}

Details: amino-acid alkali---US 10639257; US 11497698; WO 1997/004739. Low-sensitizing precursor---Goebel et al.\ 2014. Resorcinol-free---US 11517514. Melanin-mimetic---Brown et al.\ 1989; US 5346509; RSC Adv.\ 2019, 2021. Encapsulated oxidant---US 4844711; US 5053051; US 10653592.

ME-PPD is the closest prior art for a lower-sensitizing precursor and is therefore the benchmark here. It still gives approximately 30\% cross-elicitation in PPD-allergic subjects (Bl\"omeke et al.\ 2015; Kock et al.\ 2016) and was not designed for environmental degradation. Amino-acid alkalizers address odor and volatility, not precursor fate. Melanin-mimetic systems avoid aromatic amines but cover a narrow gamut and behave mainly as surface coatings.

Gap statement: ``To our knowledge, no published work designs oxidative-dye precursors and chromophores with distinct, intentional fates in hair, on skin and in wastewater. This is the focus of the present proposal.''

\section{Research question and thesis}

\textbf{Research question (RQ):} Can an oxidative dye precursor be designed to be stable in hair, esterase-labile on skin, and its coupled dye cleavable in wastewater, without losing coupling kinetics or wash-fastness relative to ME-PPD?

\textbf{Thesis:} Single-property substitutions have not resolved the performance-versus-safety trade-off. Conditional stability, meaning stable in one environment and labile in another, is the design principle that can.

\section{Objectives and hypotheses}

\textbf{O1:} Synthesize 5--8 PPD analogues with cleavable groups (ester, carbonate, hindered amide).\\
\textbf{O2:} Show coupling kinetics within $\pm$25\% of ME-PPD with m-aminophenol-type couplers [TO MEASURE].\\
\textbf{O3:} Show esterase half-life on skin-like models of $\leq$60 min and sensitization potency (DPRA, KeratinoSens, h-CLAT) no worse than ME-PPD [TO MEASURE].\\
\textbf{O4:} Show 28-shampoo fastness $\Delta E^*_{ab} \leq 3$ against an ME-PPD control on natural and white tresses.\\
\textbf{O5:} Show $\geq 60\%$ primary degradation of the coupled dye within 28 days under OECD 301B-type conditions [TO MEASURE].

\textbf{H1:} A sterically hindered ester survives 30--45 min at pH 9--10 inside hair but is cleaved by skin esterases.\\
\textbf{H2:} The cleaved carboxylate or alcohol product is more polar and less able to haptenate proteins.\\
\textbf{H3:} A trigger link placed away from the chromophore does not shift color or fastness.

\section{Design concept}

This section replaces any four-pillar system description. The contribution is the precursor and chromophore design only.

\begin{itemize}[leftmargin=1.5em]
\item \textbf{Core:} a PPD or ME-PPD scaffold that keeps the oxidation potential needed for coupling.
\item \textbf{Cleavable side chain:} an ester or carbonate at the 2-position or on an N-alkyl group, tuned by steric hindrance.
\item \textbf{Trigger link in the chromophore:} a link that is stable at pH 9 in hair but breaks under enzymatic or pH-shift conditions in effluent.
\item \textbf{Baseline system:} arginine-based alkali and conventional 6\% peroxide for the controls. Alkali and oxidant optimization is future work.
\end{itemize}

\textbf{Key risk, stated openly:} esters can hydrolyze at pH 9--10 before they reach the cortex. The screening funnel must test stability in the formulation first. Candidates that hydrolyze in the tint base before application fail early, regardless of esterase behavior.

\section{Experimental plan (12 months)}
\label{sec:plan}

\subsection{Phase 1 (months 1--4): synthesis and in-vitro}
\begin{itemize}[leftmargin=1.5em]
\item Synthesize analogues; confirm structure and purity.
\item UV-Vis / LC-MS coupling kinetics vs.\ ME-PPD.
\item Ester stability at pH 9--10 in formulation-like buffer.
\item Esterase assay on skin-like models.
\item DPRA, KeratinoSens, h-CLAT; cytotoxicity (3T3 NRU).
\end{itemize}

\subsection{Phase 2 (months 5--9): tresses}
\begin{itemize}[leftmargin=1.5em]
\item Substrates: natural, bleached, and 100\% white tresses.
\item Color yield (CIELAB), 28-shampoo fastness, light fastness.
\item Tensile strength and combing force.
\item Use $n \geq 5$ tresses per condition, with controls (ME-PPD dye, PTD dye) and statistical tests (ANOVA with post-hoc).
\end{itemize}

\subsection{Phase 3 (months 10--12): fate testing}
\begin{itemize}[leftmargin=1.5em]
\item OECD 301B on precursors and coupled dyes separately.
\item LC-MS identification of breakdown products.
\item Screening-level aquatic toxicity.
\end{itemize}

\textbf{Out of scope:} clinical HRIPT and salon trials (listed as future work), and any claim of Level 10 lift.

\section{Expected results and success criteria}

\begin{table}[h]
\centering
\caption{Success criteria vs.\ ME-PPD / PTD control}
\label{tab:criteria}
\begin{tabular}{@{}lcc@{}}
\toprule
Metric & Control & Target \\
\midrule
Coupling rate & 1.0 (ref) & 0.75--1.25 \\
28-shampoo $\Delta E^*_{ab}$ & measured & $\leq 3$ ($\leq 4$ max) \\
Grey coverage, white & measured & $\geq 95\%$ \\
Sensitization in-vitro & ME-PPD value & not worse \\
Dye degradation, 28 d & $\sim$0\% & $\geq 60\%$ \\
Tensile loss & measured & not worse \\
\bottomrule
\end{tabular}
\end{table}

Notes: control is ME-PPD / PTD dye as specified in Section~\ref{sec:plan}. Coupling rate is normalized to ME-PPD. Dye primary degradation is measured under OECD 301B-type conditions at 28 days.

\section{Limitations}

\begin{itemize}[leftmargin=1.5em]
\item Ester stability at alkaline pH may force a narrower design window.
\item In-vitro sensitization assays do not replace clinical testing.
\item New precursors need a full regulatory dossier, which is outside this scope.
\item Lift beyond 1--2 levels and platinum blonde are not targeted.
\end{itemize}

\section{IP statement}

``Related patents on amino-acid alkalizers, ME-PPD, resorcinol-free dyes and encapsulated oxidants exist. The present work aims to claim the cleavable-precursor and trigger-link chromophore concept. A freedom-to-operate review is required before commercialization.''

\section{Sustainability and safety framework}

\begin{itemize}[leftmargin=1.5em]
\item Non-animal sensitization battery only (OECD TG 442C/D/E); no new animal testing proposed.
\item Ready biodegradability assessed by OECD 301B on precursors and coupled dyes separately, not on whole formula alone.
\item LCA screening follows ISO 14040 logic (cradle-to-rinse comparison vs.\ ME-PPD control) without claiming a full ISO-compliant LCA in 12 months.
\item Green chemistry audit limited to what this scope controls: atom economy of the analogue synthesis, E-factor of the screening route, and removal of the persistent-chromophore endpoint.
\end{itemize}

\section{Risks and mitigations}

\begin{table}[h]
\centering
\caption{Risks and mitigations}
\label{tab:risks}
\begin{tabular}{@{}lcl@{}}
\toprule
Risk & Likelihood & Mitigation \\
\midrule
Ester hydrolyzes at pH 9--10 & High & Hindered ester; stability gate first \\
Coupling rate drops & Medium & Keep PPD core; 5--8 funnel \\
In-vitro sensitization equivocal & Medium & Three-assay battery + esterase $t_{1/2}$ \\
Trigger link shifts color & Medium & Link off-chromophore; CIELAB gate \\
301B pass precursor, fail dye & Medium & Test both separately; redesign trigger \\
\bottomrule
\end{tabular}
\end{table}

\section{Future work}

\begin{itemize}[leftmargin=1.5em]
\item Arginine-alkali and low-dose oxidant optimization around the lead precursor.
\item Clinical HRIPT and salon half-head evaluation after in-vitro and tress gates pass.
\item Full shade-chart extension beyond the m-aminophenol-type coupler set.
\item Full regulatory dossier and ISO 14040/14044 LCA at pilot scale.
\end{itemize}

\section{Conclusion}

Permanent oxidative dyes persist because small aromatic amines couple into wash-fast chromophores, and persist where they should not: on skin as haptens and in effluent as stable aromatics. This proposal tests whether conditional stability can separate those fates: stable in hair for 28-shampoo fastness, esterase-labile on skin, and trigger-cleavable in wastewater. The test is comparative against ME-PPD at every gate (kinetics, sensitization in-vitro, tress performance, and OECD 301B degradation), with negative outcomes defined in advance.

\begin{thebibliography}{99}
\bibitem{goebel2014} Goebel C.\ et al.\ Introduction of a methoxymethyl side chain into p-phenylenediamine attenuates its sensitizing potency.\ \textit{Toxicol Appl Pharmacol} 2014;274(3):480--487.
\bibitem{bloemeke2015} Bl\"omeke B.\ et al.\ Cross-elicitation responses to ME-PPD in PPD-allergic individuals.\ \textit{Br J Dermatol} 2015;172(4):976--980.
\bibitem{kock2016} Kock M.\ et al.\ Continuous usage of a hair dye containing ME-PPD.\ \textit{Br J Dermatol} 2016;174(5):1042--1050.
\bibitem{gargano2022} Gargano E.M.\ et al.\ ME-PPD and hairdresser sensitization risk.\ \textit{Dermatitis} 2022.
\bibitem{brown1989} Brown K.\ et al.\ Hair coloring by melanin precursors.\ \textit{J Soc Cosmet Chem} 1989;40:65--74.
\bibitem{gao2019} Gao Z.F.\ et al.\ Polydopamine coating as a multifunctional hair dye.\ \textit{RSC Adv} 2019;9:20492.
\bibitem{rsc2019b} Melanin-mimetic multicolor and low-toxicity hair dye.\ \textit{RSC Adv} 2019;9:33617--33624.
\bibitem{draelos} Draelos Z.D.\ et al.\ MEA-based ammonia-free permanent hair dye.\ \textit{J Cosmet Dermatol} (PMC9796119).
\bibitem{pat-alkali} US 10639257; US 11497698; WO 1997/004739---amino-acid alkalizers.
\bibitem{pat-res} US 11517514---resorcinol-free dyes.
\bibitem{pat-mel} US 5346509---melanin-precursor dyeing.
\bibitem{pat-ox} US 4844711; US 5053051; US 10653592---encapsulated oxidants.
\bibitem{pat-extra} US 10952946.
\bibitem{oecd} OECD TG 301B (ready biodegradability); OECD TG 442C/D/E (DPRA, KeratinoSens, h-CLAT).
\bibitem{iso14040} ISO 14040---life cycle assessment principles.
\end{thebibliography}

\end{document}
